% ---------------------------------------------------------------------------
% Section 2: the unified formulation ("Frozen Basis Time Series Forecasting").
% Implements notation_and_taxonomy_spec.md §2–§5, §7; paper_rewrite_plan.md
% §6 subplan "§2 The Method" and §13 step 3.
% Symbol closure: every symbol is defined exactly once in
% Table~\ref{tab:dimensions}; quantities outside it are either derived
% (\hat{x}, \hat{y}, block index ell) or locally defined pre-training-predictor
% notation (Sec. method_predictors). No retired symbol (sigma, gamma, Psi)
% appears anywhere in this fragment.
% ---------------------------------------------------------------------------

\section{The Method: Two Integers and a Trunk Choice}
\label{sec:method}

N-BEATS \citep{oreshkin2020nbeats} reads a time series through a doubly
residual stack of blocks. Each block receives the current residual stream
$\mathbf{x}$, emits a backcast $\hat{\mathbf{x}}$ of the input window and a
forecast $\hat{\mathbf{y}}$ of the horizon, passes $\mathbf{x} -
\hat{\mathbf{x}}$ downstream, and contributes $\hat{\mathbf{y}}$ to the model
forecast. Throughout this paper the \emph{trunk} is the multilayer mapping
inside a block that turns its input into a hidden representation, and the
\emph{basis} is a fixed matrix whose rows span everything a block is allowed
to emit. The submitted version spread the description of these blocks across
three sections; this section states the method once, before any variant is
named.

Every block in this paper computes the same three steps,
\begin{equation}
\label{eq:formulation}
\underbrace{\mathbf{h} = \operatorname{trunk}(\mathbf{x})}_{\text{trunk}};\qquad
\underbrace{\boldsymbol{\theta}_b = \mathbf{W}_b\,\mathbf{h},\;\;
\boldsymbol{\theta}_f = \mathbf{W}_f\,\mathbf{h}}_{\text{learned, unconstrained}};\qquad
\underbrace{\hat{\mathbf{x}} = \mathbf{B}_b^{\top}\boldsymbol{\theta}_b,\;\;
\hat{\mathbf{y}} = \mathbf{B}_f^{\top}\boldsymbol{\theta}_f}_{\text{frozen basis}},
\end{equation}
where the projections $\mathbf{W}_b$ and $\mathbf{W}_f$ are learned without
constraint and the basis matrices $\mathbf{B}_b \in \mathbb{R}^{d_b \times L}$
and $\mathbf{B}_f \in \mathbb{R}^{d_f \times H}$ are frozen at build time with
orthonormal rows. Blocks are grouped into $S$ stacks (default 10); writing
$(\hat{\mathbf{x}}_\ell, \hat{\mathbf{y}}_\ell)$ for the outputs of block
$\ell$, the stream updates as $\mathbf{x}_{\ell+1} = \mathbf{x}_\ell -
\hat{\mathbf{x}}_\ell$ while the model forecast accumulates the
$\hat{\mathbf{y}}_\ell$. Under \eqref{eq:formulation} the entire method
reduces to \emph{two integers and one trunk choice}: the ranks $d_b$ and
$d_f$, and whether the trunk has uniform hidden width $U$ or an autoencoder
bottleneck of latent width $d$. Every variant studied in this paper is a
choice of how the rows of $\mathbf{B}$ are generated, or of how $\mathbf{h}$
is produced. Table~\ref{tab:dimensions} defines every symbol used in this
paper; no other table or equation introduces one.

\begin{table}[t]
\centering
\small
\caption{Dimension table. Every mathematical symbol in the paper appears here
exactly once; nothing outside this table may introduce a symbol. Ranks follow
the shorthand of the results tables: $d_b$ and $d_f$ are written $H$ when
equal to the horizon, $2H$ when twice it, $L$ when matched to the input
window, and as an integer otherwise.}
\label{tab:dimensions}
\begin{tabular}{lll}
\toprule
Symbol & Space / value & Meaning \\
\midrule
$L$ & $\mathbb{N}$ & backcast (input window) length \\
$H$ & $\mathbb{N}$ & forecast (horizon) length \\
$S$ & $\mathbb{N}$ & number of stacks (default 10) \\
$U$ & $\mathbb{N}$ & uniform-trunk hidden width (default 512) \\
$d$ & $\mathbb{N}$ & AE-trunk latent width \\
$d_b$ & $\mathbb{N}$, ${\le}\,L$ & backcast basis rank (default $H$ or $2H$) \\
$d_f$ & $\mathbb{N}$, ${\le}\,H$ & forecast basis rank (complete when $d_f = H$) \\
$p$ & $\mathbb{N}$ & polynomial rank of the trend branch (default 3) \\
$r$ & $\mathbb{N}$ & learned generic-branch rank \\
$\mathbf{x}$ & $\mathbb{R}^L$ & block input (residual stream) \\
$\mathbf{h}$ & $\mathbb{R}^U$ or $\mathbb{R}^d$ & trunk output \\
$\boldsymbol{\theta}_b,\ \boldsymbol{\theta}_f$
  & $\mathbb{R}^{d_b},\ \mathbb{R}^{d_f}$ & coefficient vectors \\
$\mathbf{W}_b,\ \mathbf{W}_f$
  & $\mathbb{R}^{d_b\times U},\ \mathbb{R}^{d_f\times U}$ & learned projections \\
$\mathbf{B}_b,\ \mathbf{B}_f$
  & $\mathbb{R}^{d_b\times L},\ \mathbb{R}^{d_f\times H}$
  & frozen orthonormal basis matrices \\
$\mathbf{P}_b,\ \mathbf{P}_f$
  & $\mathbb{R}^{p\times L},\ \mathbb{R}^{p\times H}$
  & frozen polynomial (trend) basis matrices \\
$\mathbf{G}_f$ & $\mathbb{R}^{r\times H}$ & learned generic-branch basis \\
$\mathbf{v}$ & $\mathbb{R}^d$ & AE latent gate vector (appendix only) \\
$\phi$ & --- & trunk activation (ReLU unless stated) \\
\bottomrule
\end{tabular}
\end{table}

\subsection{Synthesis heads: unconstrained linear versus frozen basis}
\label{sec:method_heads}

The unconstrained linear synthesis head of standard N-BEATS is a limiting
case of \eqref{eq:formulation}, not a different architecture. Suppose
$d_f = H$. Then $\mathbf{B}_f$ is square and orthogonal, and for any
unconstrained head $\mathbf{M} \in \mathbb{R}^{H \times U}$ there is a
projection $\mathbf{W}_f = \mathbf{B}_f\,\mathbf{M}$ with
$\mathbf{B}_f^{\top}\mathbf{W}_f =
\mathbf{B}_f^{\top}\mathbf{B}_f\,\mathbf{M} = \mathbf{M}$ exactly (verified
numerically at residual $5.6\times10^{-15}$). At complete forecast rank with
linear synthesis, the learned projection absorbs the basis rotation: a block
with a frozen forecast basis represents precisely what standard N-BEATS
represents, in rotated coordinates, and the identity of the basis is inert.

Four scope conditions keep that claim honest, and every downstream use of the
result respects them:

\begin{enumerate}
\item \textbf{Expressiveness, not trainability.} The guarantee fixes what the
forecast head can represent, not what training finds; a complete basis can
still fail to train, which is the separate axis of
Section~\ref{sec:method_predictors}.
\item \textbf{Trend-branch-present null.} Every arm of the identity comparison
grid carried the paired polynomial trend branch, which absorbs the dominant
component that otherwise breaks coefficient conditioning; the null is a
trend-branch-present result.
\item \textbf{Forecast head only.} Every shipped backcast basis is truncated
($d_b = H$ or $2H$ rows against $L = 3$--$5H$), and truncated bases from
different generators span different subspaces that no rotation relates.
\item \textbf{Tiered configurations excluded.} Tiered offset schedules
truncate the forecast basis itself and therefore belong to the truncated
regime below, where measurement, not proof, carries the equivalence.
\end{enumerate}

Below complete rank the rotation lapses and truncation becomes the central
design choice. Setting $d_f < H$ confines every block's forecast to the span
of the retained rows of $\mathbf{B}_f$: the frozen basis makes rank an
explicit knob. An unconstrained head pins rank at $H$ implicitly and cannot
expose the trade-off; the frozen head turns it into one of the two integers
of the method. What to retain at a given rank is decided by the two quantities
of Section~\ref{sec:method_predictors}, computed before any training run.

\subsection{Basis construction}
\label{sec:method_basis}

All generators produce $\mathbf{B}_b$ and $\mathbf{B}_f$ at initialization
and register them as non-trainable buffers; neither training nor inference
ever evaluates a transform, only the matrix products of
\eqref{eq:formulation}. Three constructions are used.

\paragraph{Wavelet frames.} The discrete-wavelet generator applies inverse
reconstruction to unit impulses placed one at a time in each wavelet
coefficient band, yielding a raw synthesis frame over the target length; the
frame is then orthogonalized by singular-value decomposition and its
orthonormal rows form the basis. Row ordering matters under truncation: from
the same frame, the dyadic (band-ordered) arrangement and the
energy-descending arrangement induced by the orthogonalization span different
subspaces of equal rank. An earlier claim of ours that a wavelet basis finishes
behind random draws was an artifact of exactly this step --- our own
orthogonalization destroyed the dyadic ordering --- and is withdrawn.

\paragraph{Cosine and data-dependent bases.} The discrete cosine basis (DCT)
contributes fixed cosine rows, is defined at every horizon, and cannot
silently degenerate. The Karhunen--Lo\`eve basis (KL) diagonalizes the
uncentered second-moment matrix of forecast windows sampled from the training
set and orders rows by descending eigenvalue; among fixed orthonormal bases
it maximizes captured energy at every rank, which makes it the theoretically
optimal truncation for linear approximation.

\paragraph{Degenerate short-horizon behavior.} Below the minimum length its
filter requires, a wavelet generator returns a valid orthonormal basis whose
every row has support 1 --- the identity matrix. Several historical runs
labelled as wavelet configurations silently received this identity on short
horizons. The resulting disclosures are governed by
Section~\ref{sec:method_identity}; the revised code raises an error instead
of degenerating, and the corrected practitioner default is the DCT.
Stack-level offset schedules select consecutive row blocks for successive
stacks; the constant and tiered schedules appear as the optional
\emph{Offsets} column of the results tables, and the tiered schedule belongs
to the truncated regime by scope condition~4.

\subsection{Two pre-training predictors}
\label{sec:method_predictors}

Let $\mathbf{y}^{(1)},\dots,\mathbf{y}^{(N)}$ be a sample of training
forecast windows and let $\Pi_{\mathcal{S}}$ be the orthogonal projector onto
the span $\mathcal{S}$ of the retained rows of a basis. The
\textbf{captured energy} of that truncation is
\begin{equation}
\label{eq:capture}
\operatorname{CE}(\mathcal{S})
=\frac{\sum_{i=1}^{N}\bigl\lVert\Pi_{\mathcal{S}}\,\mathbf{y}^{(i)}\bigr\rVert^{2}}
      {\sum_{i=1}^{N}\lVert\mathbf{y}^{(i)}\rVert^{2}},
\end{equation}
the fraction of target energy lying inside the retained subspace. Writing
$\mathbf{b}_k$ for row $k$ of a complete orthonormal basis and $e_k$ for its
share of coefficient energy,
$e_k=\sum_{i}\langle\mathbf{b}_k,\mathbf{y}^{(i)}\rangle^{2}\,
/\sum_{i}\lVert\mathbf{y}^{(i)}\rVert^{2}$, the \textbf{effective rank} is
the participation ratio of the coefficient spectrum,
\begin{equation}
\label{eq:effrank}
\rho_{\mathrm{eff}}=\Bigl(\textstyle\sum_{k}e_k^{2}\Bigr)^{-1},
\end{equation}
which counts roughly how many coefficients carry the signal. Both quantities
are computable from training windows alone, before any gradient step, and
they answer the two questions that truncation raises: captured energy
predicts the trained accuracy of a truncated basis, and effective rank
predicts whether a plain linear synthesis head trains at all --- a head asked
to fill one dominant coordinate beside many small ones may fail regardless of
how much energy the basis captures. The evidence for both claims appears in
the theory and ablation sections.

\subsection{Composition and the canonical variants}
\label{sec:method_variants}

Optional branches extend \eqref{eq:formulation} by adding rows to the same
heads, never by changing its structure. The trend branch appends $p$ frozen
polynomial rows ($\mathbf{P}_b$, $\mathbf{P}_f$) to the respective basis
matrices; on a unified two-branch block the forecast path is
\begin{equation}
\label{eq:composition}
\hat{\mathbf{y}}
=\begin{bmatrix}\mathbf{B}_f\\ \mathbf{P}_f\end{bmatrix}^{\!\top}
 \boldsymbol{\theta}_f ,
\end{equation}
and the generic branch further appends $r$ learned rows
($\mathbf{G}_f$, forecast path only), giving a three-way additive
decomposition in which two structured bases stay frozen and only the
coefficients and the generic rows train. Unused branch rows are simply
absent. A two-branch block occupies each stack position (the unified layout),
or single-branch stacks alternate down the model (the alternating layout);
the layouts differ measurably and appear as the optional \emph{Layout} column
of the results tables.

The trunk is either a uniform-width network of four affine layers with
activation $\phi$, or an hourglass autoencoder that routes every input
through a bottleneck of width $d$
($U \rightarrow U/2 \rightarrow d \rightarrow U/2 \rightarrow U$), restoring
the working width at its output so that both trunks feed
\eqref{eq:formulation} identically. A gated variant multiplies the latent
coordinates elementwise by $\operatorname{sigm}(\mathbf{v})$; the gate is
measured to have no effect on accuracy (twelve paired comparisons, smallest
$p = 0.24$), so all accuracy tables report the ungated bottleneck and the
gate survives only in an appendix evidence table and the configuration
appendix. A variational latent head was implemented and is dropped from every
leaderboard: it is catastrophically worse on M4 ($5$--$10\times$), and one
appendix row records the finding.

Table~\ref{tab:variant_taxonomy} is the canonical variant list. Variant names
carry branch composition only; depth $S$, the ranks, the trunk, and the
layout are columns of the results tables, never parts of a name. The two
Oreshkin baselines keep their published names and their unconstrained heads:
they are the complete-rank limiting case of
Section~\ref{sec:method_heads}.

\begin{table}[t]
\centering
\small
\caption{Canonical variant taxonomy. Names encode branch composition only;
every other configuration axis ($S$, $d_b$, $d_f$, trunk, layout, offsets) is
a column of the results tables. The last column lists the historical encoded
names each variant subsumes; these strings appear nowhere else in the paper.}
\label{tab:variant_taxonomy}
\begin{tabular}{@{}lp{5.2cm}l@{}}
\toprule
Variant & Definition & Subsumes (historical names) \\
\midrule
N-BEATS-G & Oreshkin baseline, generic blocks &
  \texttt{NBEATS-G*} \\
N-BEATS-I+G & Oreshkin baseline, interpretable + generic &
  \texttt{NBEATS-IG*} \\
FB & frozen orthonormal basis replaces each synthesis head &
  \texttt{WaveletV3} blocks, \texttt{$\langle$fam$\rangle$WaveletV3*},
  all family-suffixed forms \\
FB+T & FB with a paired polynomial trend branch (rank $p$) &
  \texttt{TW*} (unified); \texttt{T+$\langle$fam$\rangle$V3}
  (alternating); \texttt{TWAE}, \texttt{TWAELG}, \texttt{TAE+...AE},
  \texttt{TALG}/\texttt{TAELG+...} \\
FB+T+G & FB+T plus a learned rank-$r$ generic branch &
  \texttt{TWG}, \texttt{TWGAE}, \texttt{TWGAELG} \\
NHiTS-FB & FB blocks under the NHiTS backbone &
  NHiTS sweep rows (transfer section only) \\
\bottomrule
\end{tabular}
\end{table}

\subsection{Notation notes}
\label{sec:method_notation}

Table~\ref{tab:notation_disambiguation} resolves every symbol collision of
the submitted version; no retired symbol survives with two meanings. The
submission's configuration strings packed family, stack count, latent width,
basis dimension, and activation flags into one token; that grammar is fully
replaced by the structured columns of the results tables, and the mapping
from every historical string is recorded in the configuration appendix.

\begin{table}[t]
\centering
\small
\caption{Symbol disambiguation relative to the submitted version.}
\label{tab:notation_disambiguation}
\begin{tabular}{@{}lll@{}}
\toprule
Old symbol & Collisions in the submission & Resolution here \\
\midrule
$\sigma$ & trunk activation; gate sigmoid; std in score tables &
  trunk activation $\phi$; gate written $\operatorname{sigm}(\cdot)$;
  dispersion shown only as ``$\pm$'' \\
$\gamma$ & gate vector; learning-rate decay factor &
  gate vector $\mathbf{v}$; decay written literally ``$\times0.5$'' \\
$k$ & summation index; basis rank; wavelet translation index &
  ranks $d_b$, $d_f$; $k$ reserved for summation indices \\
$\boldsymbol{\Psi}_b,\ \boldsymbol{\Psi}_f$
  & same object as the basis matrices, introduced separately &
  $\mathbf{B}_b$, $\mathbf{B}_f$ everywhere \\
$T$ & trend basis matrix; also a series length &
  trend bases $\mathbf{P}_b$, $\mathbf{P}_f$; lengths are $L$ and $H$ only \\
\texttt{trend\_thetas\_dim}, \texttt{ttd}, \texttt{td}
  & three spellings of one quantity &
  $p$; the API name appears once, in the hyperparameter appendix \\
\texttt{latent\_dim}, \texttt{ld}
  & two spellings &
  $d$; the API name appears once, in the hyperparameter appendix \\
\texttt{basis\_dim}, \texttt{bd}, \texttt{bdeq}, \texttt{bd2eq}
  & four spellings &
  $d_b$ (and $d_f$ where the two paths differ) \\
\bottomrule
\end{tabular}
\end{table}

\subsection{Identity-basis disclosure and marking}
\label{sec:method_identity}

Because a wavelet generator returns the identity below its minimum filter
length (Section~\ref{sec:method_basis}), some historical runs labelled as
wavelet configurations in fact used the identity matrix on their forecast
path. Four marking rules govern how later tables disclose this, uniformly and
without exception:

\begin{enumerate}
\item Every results-table cell whose (historical family, horizon) combination
received the identity basis carries the marker $\ddagger$ --- including the
main leaderboard, where the affected headline cells are disclosed in the
caption itself. The marker scopes exactly to the degenerate combinations: the
Coiflet and Symlet families on the Yearly, Quarterly, Weekly, Daily, and
Monthly horizons, the Daubechies-3 family on Yearly and Quarterly only, and
no Hourly cell, where every historical family was genuine.
\item The family-equivalence analysis is restated over genuinely distinct
bases per period (on Hourly all five were genuine), and the identity is
presented as the limiting case of the completeness result of
Section~\ref{sec:method_heads}, which it corroborates.
\item No sentence attributes an effect to a named basis family on an
identity-marked cell; such statements are rewritten as statements about the
variant at the stated rank, with $\ddagger$ carrying the disclosure.
\item The identity matrix itself ships as an appendix table,
cross-referenced from every marked cell, together with the corrected
practitioner rule that the DCT is the default basis because it is defined at
every horizon and cannot silently degenerate.
\end{enumerate}

Two further markers used in later tables --- an output-activation flag and a
note on historically gated trunks --- are defined once, together with the
column schema and the parameter-count convention, in the unified legend that
introduces the results section.
